% GENERATED FILE. Edit canonical lessons, then run npm run generate:books.
% canonical-source-hash: fnv1a64:664eb0791add6aef
% canonical-lessons: KA-C427-anupata, KA-C427-sekada, KA-C427-vargikarisu, KA-C427-parivartisu, KA-C427-samarpaka

\chapter{Numbers That Matter}
\label{ch:ka-numbers-that-matter}
\hlchaptermodality{427}

\begin{chapteropening}
\textbf{By the end of this chapter:} \emph{I can say a ratio, per cent, to classify, to convert and adequate.}

Complete the last lesson of chapter 427: I can say a ratio, per cent, to classify, to convert and adequate.
\end{chapteropening}

\section[\texorpdfstring{\kn{ಅನುಪಾತ} (anupāta)}{ಅನುಪಾತ (anupāta)}]{\kn{ಅನುಪಾತ} (anupāta) — a ratio; a proportion}

\label{lesson:KA-C427-anupata}



\begin{quote}
Before the new one: say the Kannada for a snag, then the Kannada for adverse.
\end{quote}

\subsection*{You'll want to know: \kn{ಅನುಪಾತ}}
\textbf{\kn{ಅನುಪಾತ}} — \emph{anupāta} — \textquotedblleft{}a ratio; a proportion\textquotedblright{}.

\begin{quote}
\kn{ನಮ್ಮ} \kn{ಶಾಲೆಯಲ್ಲಿ} \kn{ಶಿಕ್ಷಕ} \kn{ಮತ್ತು} \kn{ವಿದ್ಯಾರ್ಥಿಗಳ} \kn{ಅನುಪಾತ} \kn{ಒಂದಕ್ಕೆ} \kn{ಮೂವತ್ತು}.
\end{quote}

\begin{quote}
\emph{namma śāleyalli śikṣaka mattu vidyārthigaḷa anupāta ondakke mūvattu.}
\end{quote}

\begin{quote}
In our school the ratio of teachers to students is one to thirty.
\end{quote}

X \kn{ಮತ್ತು} Y \kn{ಅನುಪಾತ}, the ratio of X to Y.

\begin{cousinweb}
\textbf{\kn{ಅನುಪಾತ}} (\emph{anupāta}) is borrowed from Sanskrit \emph{anupāta}, \textquotedblleft{}following; proportion (in arithmetic)\textquotedblright{} (Monier-Williams).
\end{cousinweb}

\begin{grammarlens}[title={What you've built}]
One more word for saying what happens if.
\end{grammarlens}

\subsection*{Your turn}
\begin{itemize}
\raggedright
  \item \emph{Say it:} \emph{anupāta}
  \item \emph{Say it:} \emph{anupāta}, once more
  \item \emph{Say it:} \emph{namma śāleyalli śikṣaka mattu vidyārthigaḷa anupāta ondakke mūvattu.} — In our school the ratio of teachers to students is one to thirty.
  \item \emph{Recall:} say the Kannada for a snag, then the Kannada for adverse, then say \emph{anupāta} again
\end{itemize}

\begin{tcolorbox}[breakable,colback=teal!4,colframe=teal!35!black,title={Before you move on}]
What does \kn{ಅನುಪಾತ} mean? (\textquotedblleft{}A ratio; a proportion\textquotedblright{}.) Say it once more.
\end{tcolorbox}
\section[\texorpdfstring{\kn{ಶೇಕಡ} (śēkaḍa)}{ಶೇಕಡ (śēkaḍa)}]{\kn{ಶೇಕಡ} (śēkaḍa) — per cent}

\label{lesson:KA-C427-sekada}



\begin{quote}
Before the new one: say the Kannada for adverse, then the Kannada for a ratio.
\end{quote}

\subsection*{You'll want to know: \kn{ಶೇಕಡ}}
\textbf{\kn{ಶೇಕಡ}} — \emph{śēkaḍa} — \textquotedblleft{}per cent\textquotedblright{}.

\begin{quote}
\kn{ಈ} \kn{ವರ್ಷ} \kn{ಶೇಕಡ} \kn{ತೊಂಬತ್ತು} \kn{ವಿದ್ಯಾರ್ಥಿಗಳು} \kn{ಪಾಸಾದರು}.
\end{quote}

\begin{quote}
\emph{ī varṣa śēkaḍa tombattu vidyārthigaḷu pāsādaru.}
\end{quote}

\begin{quote}
This year ninety per cent of the students passed.
\end{quote}

It comes before the number: \kn{ಶೇಕಡ} \kn{ಹತ್ತು}, ten per cent.

\begin{cousinweb}
\textbf{\kn{ಶೇಕಡ}} (\emph{śēkaḍa}) is usually traced to a Hindustani or Marathi word for \textquotedblleft{}per hundred\textquotedblright{}, which goes back to Sanskrit \emph{śata}, \textquotedblleft{}a hundred\textquotedblright{} (CDIAL 12278).
\end{cousinweb}

\begin{grammarlens}[title={What you've built}]
One more word for saying what happens if.
\end{grammarlens}

\subsection*{Your turn}
\begin{itemize}
\raggedright
  \item \emph{Say it:} \emph{śēkaḍa}
  \item \emph{Say it:} \emph{śēkaḍa}, once more
  \item \emph{Say it:} \emph{ī varṣa śēkaḍa tombattu vidyārthigaḷu pāsādaru.} — This year ninety per cent of the students passed.
  \item \emph{Recall:} say the Kannada for adverse, then the Kannada for a ratio, then say \emph{śēkaḍa} again
\end{itemize}

\begin{tcolorbox}[breakable,colback=teal!4,colframe=teal!35!black,title={Before you move on}]
What does \kn{ಶೇಕಡ} mean? (\textquotedblleft{}Per cent\textquotedblright{}.) Say it once more.
\end{tcolorbox}
\section[\texorpdfstring{\kn{ವರ್ಗೀಕರಿಸು} (vargīkarisu)}{ವರ್ಗೀಕರಿಸು (vargīkarisu)}]{\kn{ವರ್ಗೀಕರಿಸು} (vargīkarisu) — to classify}

\label{lesson:KA-C427-vargikarisu}



\begin{quote}
Before the new one: say the Kannada for a ratio, then the Kannada for per cent.
\end{quote}

\subsection*{You'll want to know: \kn{ವರ್ಗೀಕರಿಸು}}
\textbf{\kn{ವರ್ಗೀಕರಿಸು}} — \emph{vargīkarisu} — \textquotedblleft{}to classify\textquotedblright{}.

\begin{quote}
\kn{ಗ್ರಂಥಪಾಲಕ} \kn{ಪುಸ್ತಕಗಳನ್ನು} \kn{ವಿಷಯದ} \kn{ಪ್ರಕಾರ} \kn{ವರ್ಗೀಕರಿಸಿದನು}.
\end{quote}

\begin{quote}
\emph{granthapālaka pustakagaḷannu viṣayada prakāra vargīkarisidanu.}
\end{quote}

\begin{quote}
The librarian classified the books by subject.
\end{quote}

Past (I …): \kn{ವರ್ಗೀಕರಿಸಿದೆ}, \emph{vargīkariside}.

\kn{ವರ್ಗ} (class) + -\kn{ಈಕರಿಸು} (to make into).

\begin{cousinweb}
\textbf{\kn{ವರ್ಗೀಕರಿಸು}} (\emph{vargīkarisu}) is built on Sanskrit \emph{varga}, \textquotedblleft{}a group, a class\textquotedblright{} (CDIAL 11331), with -\emph{īkarisu}, \textquotedblleft{}to make into\textquotedblright{}.
\end{cousinweb}

\begin{grammarlens}[title={What you've built}]
One more word for saying what happens if.
\end{grammarlens}

\subsection*{Your turn}
\begin{itemize}
\raggedright
  \item \emph{Say it:} \emph{vargīkarisu}
  \item \emph{Say it:} \emph{vargīkarisu}, once more
  \item \emph{Say it:} \emph{granthapālaka pustakagaḷannu viṣayada prakāra vargīkarisidanu.} — The librarian classified the books by subject.
  \item \emph{Recall:} say the Kannada for a ratio, then the Kannada for per cent, then say \emph{vargīkarisu} again
\end{itemize}

\begin{tcolorbox}[breakable,colback=teal!4,colframe=teal!35!black,title={Before you move on}]
What does \kn{ವರ್ಗೀಕರಿಸು} mean? (\textquotedblleft{}To classify\textquotedblright{}.) Say it once more.
\end{tcolorbox}
\section[\texorpdfstring{\kn{ಪರಿವರ್ತಿಸು} (parivartisu)}{ಪರಿವರ್ತಿಸು (parivartisu)}]{\kn{ಪರಿವರ್ತಿಸು} (parivartisu) — to convert; to transform}

\label{lesson:KA-C427-parivartisu}



\begin{quote}
Before the new one: say the Kannada for per cent, then the Kannada for to classify.
\end{quote}

\subsection*{You'll want to know: \kn{ಪರಿವರ್ತಿಸು}}
\textbf{\kn{ಪರಿವರ್ತಿಸು}} — \emph{parivartisu} — \textquotedblleft{}to convert; to transform\textquotedblright{}.

\begin{quote}
\kn{ನಾನು} \kn{ಡಾಲರುಗಳನ್ನು} \kn{ರೂಪಾಯಿಗಳಿಗೆ} \kn{ಪರಿವರ್ತಿಸಿದೆ}.
\end{quote}

\begin{quote}
\emph{nānu ḍālarugaḷannu rūpāyigaḷige parivartiside.}
\end{quote}

\begin{quote}
I converted the dollars into rupees.
\end{quote}

Past (I …): \kn{ಪರಿವರ್ತಿಸಿದೆ}, \emph{parivartiside}.

Into what: -\kn{ಗೆ}. \kn{ಪರಿವರ್ತನೆ}, a change.

\begin{cousinweb}
\textbf{\kn{ಪರಿವರ್ತಿಸು}} (\emph{parivartisu}) is built on Sanskrit \emph{parivartana}, \textquotedblleft{}revolution; barter, exchange\textquotedblright{} (Monier-Williams), with the Kannada verb ending -\emph{isu}.
\end{cousinweb}

\begin{grammarlens}[title={What you've built}]
One more word for saying what happens if.
\end{grammarlens}

\subsection*{Your turn}
\begin{itemize}
\raggedright
  \item \emph{Say it:} \emph{parivartisu}
  \item \emph{Say it:} \emph{parivartisu}, once more
  \item \emph{Say it:} \emph{nānu ḍālarugaḷannu rūpāyigaḷige parivartiside.} — I converted the dollars into rupees.
  \item \emph{Recall:} say the Kannada for per cent, then the Kannada for to classify, then say \emph{parivartisu} again
\end{itemize}

\begin{tcolorbox}[breakable,colback=teal!4,colframe=teal!35!black,title={Before you move on}]
What does \kn{ಪರಿವರ್ತಿಸು} mean? (\textquotedblleft{}To convert; to transform\textquotedblright{}.) Say it once more.
\end{tcolorbox}
\section[\texorpdfstring{\kn{ಸಮರ್ಪಕ} (samarpaka)}{ಸಮರ್ಪಕ (samarpaka)}]{\kn{ಸಮರ್ಪಕ} (samarpaka) — adequate; proper}

\label{lesson:KA-C427-samarpaka}



\begin{quote}
Before the new one: say the Kannada for to classify, then the Kannada for to convert.
\end{quote}

\subsection*{You'll want to know: \kn{ಸಮರ್ಪಕ}}
\textbf{\kn{ಸಮರ್ಪಕ}} — \emph{samarpaka} — \textquotedblleft{}adequate; proper\textquotedblright{}.

\begin{quote}
\kn{ಹಳ್ಳಿಯಲ್ಲಿ} \kn{ಸಮರ್ಪಕ} \kn{ಕುಡಿಯುವ} \kn{ನೀರು} \kn{ಇಲ್ಲ}.
\end{quote}

\begin{quote}
\emph{haḷḷiyalli samarpaka kuḍiyuva nīru illa.}
\end{quote}

\begin{quote}
The village has no proper drinking water.
\end{quote}

\kn{ಸಮರ್ಪಕವಾಗಿ}, properly.

\begin{cousinweb}
\textbf{\kn{ಸಮರ್ಪಕ}} (\emph{samarpaka}) is borrowed from Sanskrit \emph{samarpaka}, \textquotedblleft{}handing over completely, imparting\textquotedblright{} (Monier-Williams); Kannada uses it for \textquotedblleft{}adequate\textquotedblright{}.
\end{cousinweb}

\begin{grammarlens}[title={What you've built}]
One more word for saying what happens if.
\end{grammarlens}

\subsection*{Your turn}
\begin{itemize}
\raggedright
  \item \emph{Say it:} \emph{samarpaka}
  \item \emph{Say it:} \emph{samarpaka}, once more
  \item \emph{Say it:} \emph{haḷḷiyalli samarpaka kuḍiyuva nīru illa.} — The village has no proper drinking water.
  \item \emph{Recall:} say the Kannada for to classify, then the Kannada for to convert, then say \emph{samarpaka} again
\end{itemize}

\begin{tcolorbox}[breakable,colback=teal!4,colframe=teal!35!black,title={Before you move on}]
What does \kn{ಸಮರ್ಪಕ} mean? (\textquotedblleft{}Adequate; proper\textquotedblright{}.) Say it once more.
\end{tcolorbox}
